\begin{definition}
For a fixed actual split-outside-blocker witness surface, cut \(k<n\), cell
\(\omega\), selected \(x\in X\), paired atom families, and paired label
events, the residual-label structure is the assertion that the atom labels
are not arbitrary formal names but the concrete labels of the outside-\(B_k\)
first-witness residual surface.

It consists of the source and target atom normal forms over the exact pieces
of \(\noderef{SamplingSplitOutsideBlockerActualFirstWitnessResidualPieces}\)
in the following generated sense.  Each source atom is exactly the set of
source states realizing one Boolean truth vector on all source label events,
and each target atom is exactly the corresponding target truth-vector atom;
the residual piece itself is not left as an uninterpreted base event inside
the atom definition.

The structure carries the finite-label assertion \(\mathcal L\) is a finite
type, and it requires both exhaustion and coverage by concrete residual
labels.  It requires a selected activation label for \(x\), labels for every
embedded-neighbor priority clause in the selected complete-witness condition,
and labels for every selected retained source blocker test.  These selected
labels are paired: activation is paired with activation of \(\phi(x)\), an
embedded-neighbor comparison \(y<x\) is paired with its \(\phi\)-transported
comparison, and a source common-blocker non-killing test at \(z_i\) is paired
with the corresponding private-blocker non-killing test at
\(\beta(x,z_i)\) whenever that private blocker is defined.

For every earlier candidate \(y<x\) the structure requires a label whose
source and target events are the complements of the earlier first-witness
residual pieces.  It requires a label for the current source and target cell
predicates \(C_\omega\) and \(C'_\omega\).  For every retained noncurrent
common source blocker test \(k<i<n\), the source event is the concrete
non-killing test at \(z_i\) and the paired target event is the corresponding
private non-killing test at \(\beta(y,z_i)\).  For every already-private
target blocker test \(i<k\), the target event is the concrete private
non-killing test at \(\beta(y,z_i)\) and the paired source event is the
matching common-blocker non-killing test at \(z_i\).  Finally, the structure
requires a label for the paired global source and target boundary-complement
events appearing in the residual definition.  Every label is also classified
into one of these concrete paired families.  Thus a node assuming this
structure is assuming that the residual surface is generated by explicit
finite labels tied to the concrete source and target events, including the
global boundary complements, rather than by an arbitrary paired truth-vector
atomization with the whole residual piece hidden in the base.
\end{definition}
